\section{Deviations from Harris \& Harris}
\label{sec:textbook-deviations}

This implementation is based on the Harris \& Harris pipelined reference, with
intentional changes to keep the RTL modular, more ISA-complete for labs, and
synthesis-ready for FPGA BRAM inference.

\subsection{Summary table}
\begin{longtable}{@{}p{0.22\textwidth}p{0.72\textwidth}@{}}
\toprule
Area & Deviation (this repository) \\
\midrule
\endfirsthead
\toprule
Area & Deviation (this repository) \\
\midrule
\endhead
ISA coverage &
All six RV32I branch types are supported (including unsigned branches), and
loads/stores cover byte/half/word widths with sign/zero extension. \\
\addlinespace
Control organization &
Control is decoded in D and explicitly registered into E/M/W; signal names are
stage-suffixed throughout (\texttt{ImmSrcD}, \texttt{ALUControlE},
\texttt{MemWriteM}, \texttt{ResultSrcW}). \\
\addlinespace
Branching &
A dedicated branch unit evaluates all branch types instead of relying only on
an ALU \texttt{Zero} flag. \\
\addlinespace
Datapath flexibility &
The ALU A input supports RS1/PC/zero selection (\texttt{SrcASelE}) for clean
PC-relative and LUI-style flows. \\
\addlinespace
Structural datapath &
\texttt{datapath\_pipe.vhd} uses a structural architecture that instantiates
five per-stage wrapper entities (\texttt{fetch\_stage} through
\texttt{writeback\_stage}), each containing named \texttt{pipe\_reg} instances
and functional-unit instantiations. This gives Vivado's elaborated design view
a navigable per-stage hierarchy; the textbook keeps the datapath as a single
flat module. \\
\addlinespace
Hazards &
A dedicated hazard unit implements forwarding, load-use stalls, and control-hazard flushing; M-stage forwarding is suppressed for loads because load data is only valid in W. \\
\addlinespace
FPGA-ready memories &
Instruction and data memories use synchronous reads so the same RTL can infer
Block RAM on FPGA tools (Vivado/Quartus). Output registers at the memories act
as the relevant pipeline latches (\texttt{InstrD} and \texttt{ReadDataW}). The
instruction memory also has an explicit enable so its registered output can be
held stable across decode-stage stalls (textbook-style ``freeze InstrD''). \\
\addlinespace
Dual-port instruction memory &
\texttt{imem\_dp.vhd} provides two synchronous read ports backed by a single
memory array: port A serves instruction fetch while port B serves data-side
loads targeting the text region (for \texttt{.rodata} constants). This avoids
duplicating IMEM in synthesis. \\
\addlinespace
Flush around BRAM outputs &
Because BRAM outputs cannot be overridden like a conventional F$\rightarrow$D register, a small \texttt{flush\_d\_hold} mechanism inserts a NOP on the following cycle to implement the same logical flush behavior. \\
\addlinespace
Byte-write inference &
Stores use a byte-lane write pattern aligned with Xilinx BRAM inference templates (direct per-byte enable assignments). \\
\addlinespace
Explicit read enables &
Data memory includes a \texttt{MemRead} enable so loads are explicitly qualified rather than being read-always. \\
\addlinespace
Register file timing &
The register file writes on the rising edge (single-edge, BRAM-friendly) and
uses a combinational bypass on each read port to preserve same-cycle
W$\rightarrow$D read-after-write behavior. \\
\addlinespace
Harvard split exposure &
Instruction and data ports are split at the top-level wrapper (\texttt{top.vhd}) while the internal core wiring remains close to textbook signal naming. \\
\bottomrule
\end{longtable}

\subsection{Hazard-unit design notes (practical guards)}
The documentation calls out a few guards that are easy to miss when moving from
diagrams to working RTL:
\begin{itemize}[leftmargin=*]
  \item M-stage forwarding excludes loads (load data is only valid in W).
  \item RAW stall checks are gated by \texttt{RegWrite} to avoid false dependencies on encodings that carry a non-zero \texttt{Rd} field but do not write a register.
  \item A control-flow redirect (taken branch/jump) must override data stalls so the pipeline advances to the correct target and makes forward progress.
\end{itemize}

\subsection{True dual-port IMEM for \texttt{.rodata} loads}
The core places \texttt{.rodata} in instruction memory alongside \texttt{.text},
so some loads must read from IMEM (addresses below \texttt{DMEM\_BASE\_ADDR}) and
others from DMEM. The top-level wrapper uses \texttt{imem\_dp} (true dual-port)
so instruction fetch and a data-side read can happen simultaneously without
duplicating BRAM resources.

\paragraph{Why this matters on FPGA}
Instantiating two separate single-port IMEMs (one for fetch, one for data-side
reads) would infer two independent BRAM blocks in synthesis. A true dual-port
memory keeps one shared array:

\begin{center}
\small
\begin{tabular}{@{}p{0.58\textwidth}p{0.28\textwidth}@{}}
\toprule
Configuration & BRAM primitives (example) \\
\midrule
Two separate \texttt{imem} instances & 2$\times$ \\
One \texttt{imem\_dp} dual-port instance & 1$\times$ \\
\bottomrule
\end{tabular}
\end{center}
\normalsize

\paragraph{Scaling (IMEM + DMEM)}
At 32-bit word width on Xilinx 7-series devices, a \texttt{RAMB36E1} stores
roughly 1024 words, so BRAM usage scales approximately as
\texttt{WORDS / 1024} per memory.
\begin{itemize}[leftmargin=*]
  \item With the default sizing (\texttt{IMEM\_WORDS=2048},
        \texttt{DMEM\_WORDS=2048}), \texttt{imem\_dp} is about 2 BRAM36 blocks
        while \texttt{dmem} is about 2.
  \item Duplicating instruction memory (two separate \texttt{imem} instances)
        would double the IMEM cost (e.g., about 4 BRAM36 blocks for
        \texttt{IMEM\_WORDS=2048}).
\end{itemize}

\subsection{Register file: rising-edge write with bypass}
The Harris \& Harris reference writes the register file on the falling edge so
a same-cycle Decode-stage read sees the freshly written value. This repository
writes on the rising edge (single-edge clocking for BRAM/timing closure) and
adds a simple combinational bypass per read port for the same-cycle
W$\rightarrow$D case:

\begin{lstlisting}[style=minrv,language=VHDL]
rs1_data <= rd_data when (rd_write = '1' and rs1_addr = rd_addr and rs1_addr /= "00000")
            else read_reg(regs, rs1_addr);
\end{lstlisting}

\subsection{Textbook vs BRAM-style timing at the stage boundary}
\begin{center}
\small
\begin{tabular}{@{}p{0.44\textwidth}p{0.44\textwidth}@{}}
\toprule
Textbook (combinational) & This design (synchronous BRAM) \\
\midrule
\texttt{InstrF} available same cycle & BRAM delivers \texttt{InstrD} on the next rising edge \\
Separate F$\rightarrow$D register for \texttt{InstrD} & IMEM output register \emph{is} the F$\rightarrow$D latch \\
Separate M$\rightarrow$W register for \texttt{ReadDataW} & DMEM output register \emph{is} the M$\rightarrow$W latch \\
Read-always memory & Read enables qualify and/or freeze outputs on stalls \\
\bottomrule
\end{tabular}
\end{center}
\normalsize

\paragraph{Does this add a pipeline stage?}
No. Instead of inserting a 6th stage just to ``wait for BRAM'', the synchronous
memory output registers become the stage-boundary latches:
\begin{itemize}[leftmargin=*]
  \item IMEM: during cycle N, \texttt{PCF} is presented as the read address; on
        the next rising edge (cycle N+1), the BRAM output register updates and
        the instruction is available as \texttt{InstrD} for Decode.
  \item DMEM (loads): during cycle N, the load address in \texttt{ALUResultM} is
        presented; on the next rising edge (cycle N+1), the BRAM output register
        updates and the load word is available as \texttt{ReadDataW} for
        Writeback.
\end{itemize}
\noindent
The stage names (F/D/E/M/W) describe when values are valid and consumed, not
when a BRAM address is presented.

\subsection{Why the synchronous-memory change matters}
The textbook's combinational-read arrays are convenient for diagrams but do not
map cleanly to FPGA Block RAM primitives. This project adopts BRAM-style,
clocked reads to keep the educational pipeline \emph{real-hardware} friendly,
while preserving the same stage timing and hazard behavior.
